\documentclass[12pt]{article}
\usepackage[margin=1in]{geometry}
\usepackage{graphicx}
\usepackage{booktabs}
\usepackage{amsmath}
\usepackage{hyperref}
\usepackage[utf8]{inputenc}

\title{
Nasion Axis Consistency in Static Precision vs Dynamic Propulsion Athletes\\
\large A Constitutional Mechanics Comparative Study of Olympic Gold Medalists
}

\author{
DuckJin Yoon (Solvaram)\\
Eight-Constitution Research Platform
}

\date{2026}

\begin{document}

\maketitle

\begin{abstract}
Athletic disciplines are conventionally divided into static precision and dynamic propulsion categories based on surface biomechanics. 
This study proposes a constitutional reinterpretation: athletic expression emerges from organ dominance represented spatially by Nasion axial polarity. 
Olympic gold medalists from archery, shooting, sprint, and swimming were comparatively examined. 
Static precision disciplines consistently align with Vesicotonia (3S polarity), sprint propulsion aligns with Cholecystonia (2N polarity), and aquatic propulsion aligns with Pulmotonia (1N polarity). 
Rather than opposing static and dynamic categories, this framework demonstrates that dynamic performance is layered upon constitutional stabilization. 
The study presents structural coherence as a conceptual foundation for future single-sport axial purity investigations.
\end{abstract}

\section{Introduction}

Precision sports such as archery and shooting emphasize postural stability and fine motor control. 
Biomechanical studies confirm the relationship between body stability and shooting performance \cite{santos2025_archery_stability,fan2025_archery_postural}. 
Dynamic-static balance perspectives have also been discussed in archery research \cite{kesilmis2024_archery_balance}.
 

However, these approaches remain biomechanical rather than constitutional. 
This study proposes that athletic specialization originates from organ dominance patterns expressed through the Nasion axial coordinate system.

\section{Theoretical Framework}

Constitutional mechanics follows the structural sequence:

\begin{center}
Organ Dominance $\rightarrow$ Neuromuscular Regulation $\rightarrow$ Tension Pattern $\rightarrow$ Athletic Expression
\end{center}

The Nasion axis serves as a spatial polarity marker reflecting constitutional dominance.

\section{Comparative Axial Analysis}

\subsection{Static Precision: Vesicotonia (3S)}

Disciplines: Archery, Shooting

Functional characteristics:
\begin{itemize}
\item Sustained isometric tension
\item Pelvic stabilization dominance
\item Micro-vibration suppression
\end{itemize}

Physical modeling of archery dynamics supports the importance of stability regulation \cite{archery_physics_meyer2015}.

\subsection{Sprint Propulsion: Cholecystonia (2N)}

Discipline: 100m Sprint

Functional characteristics:
\begin{itemize}
\item Acceleration ignition
\item Gallbladder-dominant torque
\item Rapid forward propulsion
\end{itemize}

Explosive sprint initiation reflects ignition-based constitutional mechanics rather than generalized dynamism.

\subsection{Aquatic Propulsion: Pulmotonia (1N)}

Discipline: Olympic Swimming

Functional characteristics:
\begin{itemize}
\item Thoracic propulsion dominance
\item Respiratory synchronization
\item Rhythmic upper-body drive
\end{itemize}

Performance modeling in shooting-endurance integration contexts further supports layered regulation patterns \cite{leonelli2024_biathlon_bayesian}.

\section{Static-Dynamic Reinterpretation}

Static and dynamic are not oppositional constructs. 
Dynamic elevation is built upon constitutional stabilization. 
Pole vaulting, though dynamically expressive, relies on vesicotonic static dominance during pole compression phases.

Thus:

\begin{center}
Dynamic Expression = Stabilization Base + Propulsive Mode
\end{center}

\section{Implications}

\begin{itemize}
\item Constitutional prediction of sport suitability
\item Training personalization based on axial dominance
\item Injury-pattern forecasting
\end{itemize}

\section{Scope and Future Direction}

This study demonstrates structural coherence rather than statistical mass validation. 
Future research will investigate single-sport axial purity and axis-specific clustering patterns.

\section*{Project Archive}

Project materials available at:

\url{https://doi.org/10.17605/OSF.IO/9QBDP}

\bibliographystyle{plain}
\bibliography{BIMOD_200_references}

\end{document}
